\begin{table}[H]
\centering
\caption{Lista de resultados para el \autoref{sec:formulación_ej1}}
\label{tab:res_ej1}
\begin{threeparttable}
{\setstretch{1.5}
\begin{tabular}{p{0.16\linewidth}p{0.16\linewidth}p{0.185\linewidth}p{0.185\linewidth}p{0.185\linewidth}}\hline

    \textbf{Compuesto} & \textbf{Peso molar} & \textbf{g de elemento requeridos} & \textbf{g de compuesto requeridos} & \textbf{g de compuesto en exceso} \\ \hline
    
    $C_{12}H_{22}O_{11}$ &    342g/mol & 63.11gC  &            150g           &   187.5g  \\
    $(NH_{4})_{2}SO_{4}$ & 132.14g/mol & 18.49gN  &           87.26g          & LIMITANTE\tnote{1} \\
    $K_{3}PO_{4}$        & 212.27g/mol &  3.44gP  &  \textbf{23.55g}\tnote{2} &  29.44g   \\
                         &             &  1.28gK  &            2.32g          &  -------- \\
    $MgSO_{3}$           & 104.37g/mol &  1.28gS  &   \textbf{4.17g}\tnote{2} &   5.21g   \\
                         &             &  0.64gMg &            2.78g          &  -------- \\
    $CaCl_{2}$           & 110.98g/mol &  0.38gCa &            1.05g          &   1.31g   \\
    
    \hline
\end{tabular}
}
\begin{tablenotes} % ex: \tntote{1} ... \item[1]
    \item[1] No se calculó el exceso porque es el elemento limitante, y se busca que esté solo en la cantidad justa.; 
    \item[2] Se eligió el valor más alto, pues tal cantidad de compuesto suple ambos elementos.
\end{tablenotes}
\end{threeparttable}
\end{table}